\cotefile{56}
\scanpage{2}
\margin{3 p}

Situation des ``vanishing cycles'' :

\unsure{Généralités|trivialités}

Cas des Quadriques

% archive sheet number: 1
\scanpage{3}
% slanted
\margin{\uline{Vanishing} cycles : cas des quadriques}

\diagram{\begin{tikzcd}[ampersand replacement=\&]
 X \arrow[d,"f"] \\ S
\end{tikzcd}}

$f$ un schéma projectif plat

$S$ spectre d'un anneau de val. discrète strictement local

$X_0$ fibre spéciale quadrique à point singulier

$\bar X_1$ fibre géom générale, quadrique non singulière

dimension (commune) des fibres $n-1$

\unsure{$\mathrm{I}$|$[$}

\deleted{On veut étudier}

$X\deleted{\unsure{=|{-}}}\subset P^n_S$ défini par une équation quadratique \ldots.

$H^i(X_0)\to H^i(\bar X_1)$ ? On a la suite exacte générale

\[
\cdots\to H^i(X_0)\to H^i(\bar X_1)\unsure{\mathord{-}|\to}\Phi^{i+1}\to H^{i+1}(X_0)\to H^{i+1}(\bar X_1)\to\cdots
\]

\deleted{1\textsuperscript{o}). \uline{$n$ pair}, d'où \unsure{$n-1$ impair|$n$ impair}}

On sait que la cohomologie de \deleted{$H^i(X_0)$ et} $X_0$ et $\bar X_1$ est nulle en dims impaires. Donc on trouve

\[
0\to\Phi^{2\nu}\to H^{2\nu}(X_0)\to H^{2\nu}(\bar X_1)\to\Phi^{2\nu+1}\to0
\]

i.e. les $\Phi^{i}$ apparaissent \deleted{\unsure{certes|comme}} comme les noyaux et conoyaux des \uline{hom de spécialisation}.

Notons aussi que sur $X$ donc $X_0$, on a une classe $\xi\in H^2(X,\mu(1))$, provenant de la section à l'infini, et dont la restriction aux deux fibres est connue \ldots

% archive sheet number: 2
\scanpage{4}
\mspage{\handwritten{\unsure{6 b|68|6 bis}}}

\deleted{a) La fonction $\handwritten{\rho}_r$ sur $S$ est semi-continue inférieurement, la fonction $\handwritten{\rho}_n$ sur $S$ est semi-continue supérieurement, donc la fonction $\handwritten{\rho}_u=\handwritten{\rho}_n-\handwritten{\rho}_r$ est semi-continue supérieurement.}

\deleted{b) Les \deleted{\unsure{c|e}}conditions suivantes (stables par changement de base quelconque) sont équivalentes:}

\deleted{(i) La fonction $\handwritten{\rho}_r$ sur $S$ est localement constante.}

\deleted{(ii) Il existe localement, au sens de la topologie étale, un tore maximal dans $G$.}

\deleted{(ii bis) Il existe localement, au sens de la topologie fidèlement plate quasi-compacte, un tore maximal dans $G$.}

\deleted{c) Soient $T_1,T_2$ deux tores maximaux dans $G$ (ce qui implique que l'on est sous les conditions envisagées dans b). Alors $T_1,T_2$ sont conjugués localement au sens de la topologie étale, i.e. il existe un morphisme étale surjectif $S'\handwritten{\to}S$ tel que les sous-groupes $T_{1\,S'}$ et $T_{2\,S'}$ de $G_{S'}$ soient conjugués par une section de $G_{S'}$ sur $S'$.}
% Handwritten transposition loop under "localement" and over "conjugués".

\deleted{d) Sous les conditions de b), i.e. lorsque $\handwritten{\rho}_r$ est localement constante, il en est de même de $\handwritten{\rho}_n$ (donc aussi de $\handwritten{\rho}_u=\handwritten{\rho}_n-\handwritten{\rho}_r$).}

\deleted{Démonstration.}

\deleted{a) Notons que pour tout morphisme $S'\handwritten{\to}S$, si $G'=G\mathbin{\mathrm{x}}_S S'$, les fonctions $\handwritten{\rho}'_r$ etc sur $S'$ \deleted{construites à l'aide} définies en termes de $G'$ comme $\handwritten{\rho}_r$ etc en termes de $G$, s'obtiennent simplement \deleted{\unsure{comme}} en composant ces dernières avec $S'\handwritten{\to}S$. Lorsque $S'\handwritten{\to}S$ est fidèlement plat quasi compact, il en résulte que $\handwritten{\rho}'$ est continue, resp. semi-continue supérieurement, resp. semi-continue inférieurement, si et seulement si $\handwritten{\rho}$ l'est, \deleted{(SGA VIII} la topologie \inserted{de Zariski} de $S$ étant en effet quotient de celle de $S'$ (SGA VIII 4.3.). Par suite, les assertions de a) sont locales pour}

% archive sheet number: 3
\scanpage{5}
1\textsuperscript{o}) \uline{$n$ pair} $=2\nu$, \uline{$n-1$ impair} \qquad $n-1=2\nu-1=2(\nu-1)+1$

\diagram{$\left.\begin{array}{ll}
 H^*(X_0):&1,\xi_0,\ldots,\xi_0^{\nu-1},[\xi_0^\nu,\lambda'_\nu],\frac12\xi_0^{\nu+1},\ldots\frac12\xi_0^{2\nu-1}\\[4pt]
 H^*(\bar X_1):&1,\xi_1\ldots,\xi_1^{\nu-1},\frac12\xi_1^\nu,\frac12\xi_1^{\nu+1},\ldots,\frac12\xi_1^{2\nu-1}
 \end{array}\right]\quad\begin{array}{l}\xi_0\leadsto\xi_1\\ \lambda'_\nu\leadsto0\end{array}$}

\[
\begin{aligned}
\text{Donc}\quad & H^i(X_0)\to H^i(\bar X_1)\quad\text{\uline{bijectif} \deleted{excepté} si }i\ne2\nu=n\\
& H^{2\nu}(X_0)\to H^{2\nu}(\bar X_1)\quad\text{\uline{surjectif}, noyau engendré par }\lambda'_\nu
\end{aligned}
\]

\[
\text{Donc}\quad\text{a)}\quad
\left\{\begin{aligned}
\Phi^{i}&=0\qquad \text{\deleted{sauf} si }i\ne(2\nu=n)\\
\Phi^{\unsure{2\nu|n}}&\simeq\mu(\unsure{-\nu})\qquad\text{(isom déterminé au signe près\ldots.)}
\end{aligned}\right.
\]

\[
\text{b)}\quad\left\{\begin{aligned}
&R^i f_*\quad\text{constant si }i\ne n\\
&R^n f_*\quad\text{a un sous-faisceau concentré en $0$, isom. à $\mu(\ldots)$.}\\
&\hspace{6em}\text{le quotient par ce dernier est constant}
\end{aligned}\right.
\]

\margin{c) $\pi$ opère trivialement sur $H^*(\bar X_1)$ et sur $\Phi^{2\nu}$}

2\textsuperscript{o}) $n$ impair $=2\nu+1$, \qquad $n-1=2\nu$

\diagram{$\left.\begin{array}{ll}
H^*(X_0):&1,\xi_0,\ \ ,\xi_0^{\nu-1},\xi_0^\nu,\frac12\xi_0^{\nu+1},\ldots\frac12\xi_0^{2\nu}\\[4pt]
H^*(\bar X_1):&1,\xi_1,\ldots\xi_1^{\nu-1},[\xi_1^\nu,\lambda_\nu],\frac12\unsure{\xi_1|\{_1}^{\nu+1},\ldots,\frac12\xi_1^{2\nu}
\end{array}\right]\quad\xi_0\leadsto\xi_1$}

\[
\begin{aligned}
\text{Donc}\quad & H^i(X_0)\to H^i(\bar X_1)\quad\text{bijectif si }i\ne2\nu=n-1\\
& H^{2\nu}(X_0)\to H^{2\nu}(\bar X_1)\quad\text{injectif, conoyau engendré par }\deleted{\unsure{\xi^\nu|\lambda'_\nu}}\inserted{\unsure{\ell_\nu|\lambda_\nu}}
\end{aligned}
\]

\[
\text{a)}\quad\left\{\begin{aligned}\Phi^i&=0\quad\text{si }i\ne2\nu+1=n\\
\Phi^n&\simeq\mu(?\ldots)\end{aligned}\right.
\]

\[
\text{b)}\quad\left\{\begin{aligned}
&R^i f_*\ \text{constant si }i\ne n-1\\
&R^{n-1}f_*\ \text{est ``pur'' \ldots.}
\end{aligned}\right.
\]

\margin{c) $\pi$ opère trivialement sur les $H^i(\bar X_1)$, pour $i\ne n-1$, mais non sur $\Phi^n$, \ldots.}

% archive sheet number: 4
\scanpage{6}
\mspage{\handwritten{\deleted{\unsure{II|I}} 1}}

\deleted{\handwritten{1}. \uline{Tores maximaux\handwritten{.} \deleted{et groupe de Weyl}}\handwritten{.}}

\deleted{A partir de maintenant, nous ferons usage des résultats du Séminaire Chevalley 56/57 sur les Groupes Algébriques. Suivant l'usage en théorie des groupes algébriques, nous référons à ce Séminaire \handwritten{p}ar le sigle BIBLE .}

\deleted{Soit d'abord $G$ un groupe algébrique sur un corps \uline{algébriquement clos}\deleted{\unsure{,}} $k$. On appelle \uline{tore maximal} de $G$ un \deleted{tore} sous-groupe algébrique $T$ de $G$ qui est un \uline{tore} (ce qui signifie ici, $k$ étant algébriquement clos, qu'il est isomorphe à un groupe de la forme $\underline{G}_m^{\,r}$), et qui est maximal pour cette propriété. Noter que, $k$ étant parfait, $G_{\text{réd}}$ est un sous-groupe de $G$, lisse sur $k$, c'est donc essentiellement un groupe algébrique au sens de BIBLE. D'ailleurs, tout sous-groupe réduit de $G$ (tel un tore) est automatiquement un sous-groupe de $\deleted{\unsure{G|R}}G_{\text{réd}}$, par suite les tores maximaux de $G$ coïncident avec les tores maximaux de $G_{\text{réd}}$. Lorsque $G$ est affine, donc $G_{\text{réd}}$ affine, un théorème fondamental de Borel nous dit que deux tores maximaux de $G$ sont conjugués \deleted{\unsure{sous G(|sous $G$}} par un élément de $G(k)$ \inserted{\handwritten{$=G_{\text{réd}}(k)$}} (BIBLE 6, th.4 c ), en particulier ils ont même dimension. \inserted{Nous appellerons} \deleted{L}\handwritten{l}a dimension commune des tores maximaux de $G$ \deleted{est alors appelée} le \uline{rang réductif} de $G$.}

\margin{\handwritten{$G_{\text{réd}}(k)$}}

\deleted{\deleted{\unsure{(}}Noter d'ailleurs que la restriction $G$ affine est inutile, comme il résulte d'un théorème \deleted{\unsure{de|bien}} connu de Chevalley, affirmant que tout groupe algébrique connexe lisse sur $k$ est une extension d'une variété abélienne par un groupe affine\deleted{\unsure{,}}; \deleted{nous n'aurons pas besoin de cette \unsure{généralisation}),} \inserted{\handwritten{cf N\textsuperscript{o}6 \unsure{3|1}.}}}

% slanted
\margin{\handwritten{pas à la ligne}}

\deleted{Dans \deleted{\unsure{la suite}} \inserted{\handwritten{les N\textsuperscript{os} 1 à 4}}, nous nous bornerons \deleted{aux \unsure{groupes groupes|groupes-groupes|groupes}} \inserted{\handwritten{le plus souvent aux}} préschémas en groupes affines sur la base.}

% slanted
\margin{\handwritten{à la ligne}}

\deleted{\deleted{$G$ sur $k$ étant comme ci-dessus,} \inserted{\handwritten{Soit $G$ un groupe algébrique \deleted{\illegible{1}} lisse sur $k$,}} et $T$ un tore maximal de $G$, le centralisateur \inserted{$C(T)$} de $T$ dans $G$ sera appelé le}

% archive sheet number: 5
\scanpage{7}
\uline{Calculs affines} \qquad Nous remplaçons $X$ par $U=X-Z$ (\deleted{\unsure{section}} à l'infini), on aura

\[
\cdots\to H^i(U_0)\to H^i(\bar U_1)\unsure{\mathord{-}|\to}\Phi^{i+1}\to H^{i+1}(X_0)\to H^{i+1}(\bar X_1)\to\cdots
\]

Or on a maintenant \uline{$H^i(U_0)=0$} \qquad $i\unsure{\ne|=}0$,

\[
\text{d'où}\qquad H^i(\bar U_1)\simeq\Phi^{i+1}
\]

d'où directement
\[
\begin{aligned}
\Phi^i&=0\qquad\text{si }i\ne n\\
\Phi^n&=H^{n-1}(\bar U_1)\qquad\text{engendré par le}\\
&&\text{``vanishing\unsure{,|\null} cycle'' \ldots.}
\end{aligned}
\]

% archive sheet number: 6
\scanpage{8}
\mspage{\handwritten{0}}

XII \quad \uline{Tores maximaux}, \uline{groupe de Weyl}, \uline{sous-groupes de Cartan}, \uline{centre réductif des schémas en groupes lisses et affines}.

par A. Grothendieck .

A partir du présent exposé, contrairement aux exposés précédents, nous ferons usage \deleted{\unsure{ce|le}}\handwritten{de}s résultats connus sur la structure des groupes algébriques lisses sur un corps algébriquement clos $k$, et surtout de la théorie de B\deleted{orel}\handwritten{OREL} des groupes algébriques affines, exposée dans le Séminaire Ch\deleted{evalley}\handwritten{EVALLEY} 56/57 : \deleted{sur les Groupes} ``Classification des groupes de Lie Algébriques'' . Suivant l'usage en théorie des groupes algébriques, nous référons à ce Séminaire par le sigle BIBLE. Pour les \deleted{moment,} \inserted{prochains exposés,} nous aurons besoin \inserted{notamment} des résultats de\deleted{s} BIBLE 4 , 5, 6, 7 (le chiffre après BIBLE renvoi\unsure{\handwritten{é}|e} au numéro de l'exposé).

Il \deleted{\unsure{me|ne}} semble \deleted{\unsure{pas}} \inserted{d'ailleurs} que la théorie des schémas \inserted{n'}apporte aucune simplification notable à la théorie \deleted{exposée} de B\deleted{orel}\handwritten{OREL}\deleted{\unsure{, tel qu'il|, a-t-il}} telle qu'elle est exposée dans BIBLE\handwritten{,} \deleted{C}\handwritten{c}'est pourquoi il n'a pas semblé utile de la reproduire \deleted{ici, notre objet étant} dans le présent Séminaire, notre but étant ici de déduire, de résultats connus sur les corps algébriquement clos, des résultats analogues valables sur tout préschéma de base. (Il\unsure{\,\textperiodcentered\,| }n'en sera plus de même pour la théorie de structure de C\deleted{hevalley}\handwritten{HEVALLEY}\deleted{\unsure{,}} des groupes semi-simples, qui\handwritten{,} semble-t-il\handwritten{,} se traite avec avantage ab ovo sur un préschéma de base quelconque).

\handwritten{Dans les exposés oraux \inserted{(que nous avons suivi dans les N\textsuperscript{os} 1 à 4)}, nous nous étions limités aux préschémas en groupes affines sur $S$, en nous appuyant de façon essentielle sur les théorèmes de représentabilité de Exp $\overline{\underline{\mathrm{XI}}}$ N\textsuperscript{o}4. Dans les présentes notes\deleted{\unsure{,}} (cf N\textsuperscript{os} 6 à 8) nous montrons rapidement comment on peut éliminer les hypothèses affines \deleted{\unsure{par une méthode|par une voie} \unsure{\illegible{2}|Exp XI}} \inserted{par une \deleted{\unsure{autre}} méthode} plus simple n'utilisant pas les résultats les plus délicats de Exp. $\overline{\underline{\mathrm{XI}}}$.}

% margin note written sideways
\inserted{\handwritten{\deleted{\unsure{Pour}}Pour d'autres généralisations des résultats contenus dans le présent exposé, voir également exposés \deleted{Exp.} $\overline{\underline{\mathrm{XV}}}$ et $\overline{\underline{\mathrm{XVI}}}$. (Il est évident que le contenu des exposés $\overline{\underline{\mathrm{XI}}}$, $\overline{\underline{\mathrm{XII}}}$, $\overline{\underline{\mathrm{XV}}}$, $\overline{\underline{\mathrm{XVI}}}$ devrait être complètement refondu).}}

% archive sheet number: 7
\scanpage{9}
\margin{12 p}

Calculs de cohomologie des quadriques, des cônes projetants

% archive sheet number: 8
\scanpage{11}
\uline{Cohomologie d'un cône projectif}.

\diagram{\begin{tikzcd}[ampersand replacement=\&,column sep=large]
 X \arrow[d,"h"] \& \& \\
 \overset{\deleted{\unsure{a|X}}}{C'} \arrow[d] \arrow[r,hook,dash,"i'"'] \& \hat C' \arrow[d] \arrow[r,bend left=22] \& X' \arrow[l,hook',"j'"] \arrow[d,"\unsure{\wr|?|h}"]\\
 C \arrow[r,hook,"i"'] \arrow[rr,bend right=28] \& \hat C \& X \arrow[l,hook',"j"]
\end{tikzcd}}

\[
\begin{aligned}
& i\geq0\\
&\deleted{H^{i-1}(X)\to H^i(\hat C)\to H^i(\hat C')\to H^i(X)\to H^{i+1}(\hat C)}\\
&\deleted{H^i(C)}
\end{aligned}
\]

\diagram{\begin{tikzcd}[ampersand replacement=\&,column sep=small,font=\small]
\cdots \arrow[r,dash,"\mathrm{surj}"] \& H^{i-1}(X)\arrow[r,dash]\arrow[d] \& H^i(\hat C)\arrow[r,dash,"\deleted{\unsure{\sim}}"]\arrow[d] \& H^i(\hat C')\arrow[r,"\mathrm{surj}"]\arrow[d] \& H^i(X)\arrow[r]\arrow[d,phantom,"\shortparallel"] \& \cdots\\
\& H^{i-1}(X)\arrow[r] \& H^i(C)\arrow[r,dash] \& H^i(C')\arrow[r,"\mathrm{bij}"] \& H^i(X) \&
\end{tikzcd}}

\[
\begin{aligned}E&=C-a\\ \hat E&=\hat C-a\end{aligned}
\]

\[
\boxed{H^i(C)=0\qquad i\ne0}
\]

\[
0\to H^i(\hat C)\to H^i(\hat C')\to H^i(X)\to0
\]

\[
\boxed{H^i(\hat C)\xleftarrow{\sim}H^{i-2}(X)(-1)\qquad i\ne0}
\]

\inserted{\unsure{provient|provenant} de l'injection $X\xrightarrow{j}\hat C$}

$j_*(a)j_{\unsure{a|*|2}}(b)=j_*(ab\xi)$ car $j^*j_*(a)=a\xi$, cela détermine la structure multiplicative \ldots.

\deleted{\unsure{$E$|$k$}} $\underset{\substack{\wr\\\hat C'-a}}{\hat C-a}$ est fibré sur $X$ de fibres des droites \deleted{\unsure{affi|aff.}} \inserted{vectorielles}, donc

\[
\begin{aligned}\hat C-a&=\hat E\\ \hat C-a-X&=E\end{aligned}
\]

\[
\boxed{H^*(X)\xrightarrow{\sim}H^*(X,\underbrace{\hat C-a}_{\hat E})}
\]

\diagram{\begin{tikzcd}[ampersand replacement=\&,column sep=small,font=\small]
H^{i-2}(X)(-1)\arrow[r] \& H^{\deleted{\unsure{2}}i}(\hat E)\deleted{\unsure{(E)}}\arrow[r] \& \boxed{H^i(E)\deleted{\unsure{(E)}}}\arrow[r] \& H^{i-1}(X)(-1)\arrow[r] \& H^{i+1}(\hat E)\\
\& H^i(X)\arrow[u,phantom,"\wr"] \& \& \& H^{i\unsure{+|-}1}(X)\arrow[u,phantom,"\wr"]
\end{tikzcd}}

% slanted
\margin{multiplication par $\xi$ (ou $-\xi$)}

% archive sheet number: 10
\scanpage{12}
\uline{Remarques} 1.11. a)Le préschéma $\handwritten{\mathscr T}$ de 6.10. s'appelera comme de juste le \uline{préschéma des tores maximaux} de $G$. On verra au N\textsuperscript{o}5 qu'il est en fait \uline{affine} sur $S$. On peut le vérifier directement, lorsque $G$ est isomorphe à un sous-préschéma en groupes fermé d'un préschéma de la forme $\mathrm{Gl}(n)$, en utilisant Exp XI 4.6. et en remarquant que par le raisonnement de la démonstration de 1.7. b) , $\handwritten{\mathscr T}$ s'identifie à un sous-préschéma à la fois \uline{ouvert et fermé} du préschéma $F$ qui représente les sous-groupes de type multiplicatif de $G$.

% A proofreader's slash crosses the space between des and tores in the underlined expression above.

b) Il \deleted{ne serait pas difficile} \inserted{\handwritten{est possible}} de donner une démonstration de 1.7. donc de 1.9. n'utilisant pas les résultats délicats de Exp XI, mais seulement \inserted{l\handwritten{es} \deleted{\unsure{faciles}} \handwritten{r}ésultats faci\handwritten{le}s \deleted{\unsure{de Exp XI N\textsuperscript{o}6}}} Exp XI 3.12. \inserted{\handwritten{et 6.2.,}} en travaillant uniquement avec des groupes de type multiplicatif finis sur la base (moralement, les groupes ${}_nT$ où $T$ est un tore maximal) , comparer \deleted{Exp} N\textsuperscript{o}7 . \deleted{Par contre,} \inserted{\handwritten{La même remarque vaut pour}} la démonstration de 1.10.\handwritten{,}

\margin{\handwritten{et 6.2.,}}

\deleted{s'appuie de façon essentielle sur Exp XI 4.1. Dans un autre exposé, nous donnerons une version généralisée de 1.10. par une démonstration différente (s'appuyant sur un théorème de représentabilité de EGA \unsure{V)|VI)}.}

\deleted{\deleted{Exception faite du \unsure{§, N\textsuperscript{o}5|N\textsuperscript{o}5}}, \inserted{\handwritten{De façon générale,}} \inserted{pour les résultats du N\textsuperscript{o}5,} les théorèmes de représentabilité de Exp XI N\textsuperscript{o} 4 \deleted{\unsure{me semblent}} pourrai\deleted{\unsure{ent}}\handwritten{t}\deleted{\unsure{on}} être éliminés de la théorie présentée \inserted{\handwritten{dans le présent exposé et les suivants}}}

\deleted{\handwritten{de façon générale, il est possible (et sans doute avantageux) d'éliminer totalement les résultats de représentabilité de Exp $\overline{\underline{\mathrm{XI}}}$ N\textsuperscript{o} 4 de la théorie présentée dans le présent exposé et les suivants.}}

% archive sheet number: 11
\scanpage{13}
Si $X$ est non singulier \inserted{connexe de dim $n$}, on trouve -

\[
\boxed{\begin{aligned}
 i\leq n\quad&H^i(E)\simeq P^i(X)\qquad\text{mod torsion}\\
 i\geq n+1\quad&H^i(E)\simeq \deleted{\unsure{P|R}}\unsure{R|P}^{i-1}(X)(-1)\qquad\text{mod torsion}
\end{aligned}\quad\bigg|\quad\begin{gathered}\text{\unsure{1ère|1er} version du}\\\text{th. de Lefschetz}\end{gathered}}
\]

\diagram{\begin{tikzcd}[ampersand replacement=\&,column sep=2pt,font=\scriptsize]
\cdots\arrow[r]\& \underset{\substack{\shortparallel\\0\\\text{si }i-1\ne0}}{H^{i-1}(C)}\arrow[r]\& H^{i-1}(E)\arrow[r]\& H^i_x(C)\arrow[r]\&\underset{\substack{\shortparallel\\0\\\text{si }i\unsure{=|\ne}0}}{H^i(C)}\arrow[r,dash]\&\underset{\unsure{?|R}}{H^i(E)}\arrow[r]\& H^{i+1}_x(C)\arrow[r]\&\cdots
\end{tikzcd}}

\deleted{$i\ne0,1$}

\diagram{\begin{tikzcd}[ampersand replacement=\&,column sep=2pt,font=\scriptsize]
\cdot\arrow[r]\&\underset{\substack{\shortparallel\\0}}{H^0_x(C)}\arrow[r]\& H^0(C)\arrow[r,swap,"{\substack{\mathrm{bij}\\\text{\uline{si $X$ connexe}}}}"]\&H^0(E)\arrow[r]\&H^1_x(C)\arrow[r]\&\underset{\substack{\shortparallel\\0}}{H^1(C)}\arrow[r]\&H^1(E)\arrow[r,dash]\&H^{\unsure{1|2}}_x(C)\arrow[r]\&\underset{\substack{\shortparallel\\0}}{H^{\unsure{1|2}}(C)}\\
\&\&\&\unsure{0}\arrow[ur]\&\&\&\&\&
\end{tikzcd}}

\[
\boxed{\begin{aligned}
 H^0_x(C)&=H^1_x(C)=0\\
 H^i_x(C)&\simeq H^{i-1}(E)\qquad\text{si }i\geq2
\end{aligned}}
\]

Cas $X$ une quadrique $Q^n$, donc $\hat C=Q^{\prime n+1}$, on trouve

{\small
\[
\boxed{\begin{aligned}
&\underline{H^i(Q^{\prime n+1})\xleftarrow{\sim}H^{i-2}(Q^n)(-1)}\qquad \text{si }i\ne0\qquad\text{donc: \unsure{si}}\\
&\left(1,\xi,\ldots,\xi^\nu,\frac12\xi^{\nu+1},\ldots,\frac12\xi^{2\nu}\right)\qquad\text{si }\boxed{n+1=2\nu}\\
&\left(1,\xi,\ldots,\xi^\nu,[\xi^{\nu+1},\lambda'_{\nu+1}],\frac12\xi^{\nu+2},\ldots,\frac12\xi^{2\nu+1}\right)\!:\quad\text{si }\boxed{n+1=2\nu+1}\\
&\hfill\text{avec }\left\{\begin{aligned}\xi\lambda'_{\nu+1}&=0\\{\lambda'}_{\nu+1}^{\,2}&=0\end{aligned}\right.
\end{aligned}}
\]
}

% archive sheet number: 12
\scanpage{14}
\mspage{\handwritten{\unsure{8 b|84|8}}}

dans $G$. Si $C$ est son centralisateur, alors $C$ est représentable et est lisse sur $S$ par \inserted{\handwritten{Exp. $\overline{\underline{\mathrm{XI}}}$}} 5.3. , donc la fonction $s\handwritten{\leadsto}\handwritten{\rho}_n(s)=\dim C_s$ est bien localement constante.

\deleted{Corollaire} La démonstration de \deleted{6}\handwritten{1}.7. est achevée. Nous référerons aux condition envisagées dans \deleted{6}\handwritten{1}.7. b) en disant que dans ce cas, $G$ est de \uline{rang réductif localement constant}.\deleted{\unsure{sur $S$}} Notons:

\uline{Corollaire} \deleted{6}\handwritten{1}.9. Soit $G$ comme dans \deleted{6}\handwritten{1}.7. , et soit $s\handwritten{\in}S$ tel que \inserted{$\handwritten{\rho}_u(s)=0$ i.e.} $\handwritten{\rho}_r(s)=\handwritten{\rho}_{\unsure{n|m}}(s)$,(i.e. les sous-groupes de Cartan de $G_k$ sont des tores, où $k$ est une clôture algébrique de $k(s)$). Alors il existe un voisinage \deleted{\unsure{ouver|ouvert}} ouvert $U$ de $s$ sur lequel $\handwritten{\rho}_r$ et $\handwritten{\rho}_n$ sont constants, en particulier \handwritten{p}our tout $t\handwritten{\in}U$, le rang unipotent $\handwritten{\rho}_u(t)$ de $G_t$ est nul.

En effet, \deleted{\unsure{l'hypothèse s'écrit aussi} $\unsure{\rho}_u(s)=0$, et \unsure{comme}} cela résulte immédiatement de \deleted{6}\handwritten{1}.7. a) et de l'inégalité $\handwritten{\rho}_r(t)\handwritten{\leq}\handwritten{\rho}_{\unsure{n|m}}(t)$ pour tout $t\handwritten{\in}S$.

Notons aussi que nous avons prouvé, en même temps que b), le

\uline{Corollaire} \deleted{6}\handwritten{1}.10. Soit $G$ comme dans \deleted{6}\handwritten{1}.7. , et supposons $G$ de rang réductif localement constant. Considérons le foncteur

\[
\handwritten{\mathscr T}:(\mathrm{Sch})^{\unsure{\mathrm{o}|0|\circ}}\handwritten{\longrightarrow}(\mathrm{Ens})\ ,
\]

tel que pour tout $S'$ sur $S$, on ait

\[
\handwritten{\mathscr T}(S')=\text{ensemble des tores maximaux de }G_{S'}\ .
\]

Alors $\handwritten{\mathscr T}$ est représentable par un préschéma lisse, séparé et de type fini sur $S$.

Il reste à vérifier que $\handwritten{\mathscr T}$ est \deleted{\unsure{séparé}} de type fini sur $S$\handwritten{\unsure{.|?}} La question étant locale pour la topologie fidèlement plate quasi-compacte, on peut supposer que $G$ admet un tore maximal $T$. Mais alors $\handwritten{\mathscr T}$ \deleted{\unsure{se|est}} est \deleted{\unsure{égal à|égal}} \inserted{isomorphe à} $G/\underline{\mathrm{Norm}}_G(T)$ (comparer \inserted{\handwritten{Exp $\overline{\underline{\mathrm{XI}}}$}} 5.3. bis et 5.5. )\handwritten{.} \deleted{\unsure{enfin,|en effet,}} \handwritten{L}e morphisme $G\handwritten{\longrightarrow}\handwritten{\mathscr T}$ défini par $g\handwritten{\leadsto}\mathrm{int}(g)(T)$ étant surjectif , et $G$ quasi-compact sur $S$, il en est de même de $\handwritten{\mathscr T}$ , ce qui achève la démonstration.

% archive sheet number: 13
